\section{Discussion}
The results support a qualified answer to the research questions. Students
report frequent GenAI use across development activities, but their perception
families do not form a single trajectory. Repository activity and clean
source/test change concentrate near final integration, while transcript
markers become denser in the observed 2025.2 corpus. RQ3 adds variation in
upfront architectural-artifact scope, but no robust monotonic association with
late clean-path rework or independently scored outcomes. The central
implication is not that larger artifacts guarantee better projects; it is that
AI-assisted development makes shared intent and coordination more consequential
while remaining difficult to identify causally in a small observational study.

\subsection{From local acceleration to coordination}

RQ1 and RQ2 describe two sides of AI-assisted development. Students use GenAI
for concrete tasks, yet repository-visible work intensifies near final
integration. This weakens a simple productivity story in which faster local
code production automatically yields smoother system delivery. A more
defensible interpretation is a shifted bottleneck: GenAI may reduce the effort
of producing local fragments, but it does not settle architecture, interfaces,
task boundaries, dependencies, or acceptance conditions. Those agreements
remain coordination objects in multi-person projects.

The robustness analyses also prevent reducing the late peak to one
``student syndrome'' mechanism. Some team-semesters combine high final
concentration with score gains, others show late concentration without
evaluator-score gain, and others improve with more distributed activity. The
point is not to identify one true mechanism, but to show that repository
timing, planning visibility, and evaluator trajectories separate profiles that
a simple procrastination narrative would conflate.

\subsection{Planning artifacts, rework, and complexity}

RQ3 qualifies strong planning narratives. Architectural Artifact Scope (M6a)
captures repository-visible scope, and Planning Content Evidence (M6b)
extracts structured evidence from T1 commit subjects; neither measures
architectural quality. Repository absence cannot prove that a team did not
plan elsewhere, and artifact volume cannot prove that its design was coherent
or useful. The M9 associations are therefore best read as exploratory
diagnostics: with only 14 team-semesters, coefficients are modest, unstable,
and sensitive to individual cases.

Technical complexity further bounds interpretation. Late clean-path rework can
reflect project difficulty and structure as well as planning gaps or
AI-assisted implementation dynamics. The complexity profiles do not provide a
causal adjustment, but they make explicit why rework should not be attributed
to a single source. Likewise, Clean-Path Rework (M8) follows file provenance;
it does not determine whether a change was semantically defective or
pedagogically productive.

\subsection{Why SDD is a plausible response}

The results motivate SDD as a hypothesis for future evaluation, not as a tested
treatment. Its specification-to-plan-to-task workflow matches the problem
surfaced here: when GenAI makes implementation cheaper while system-level
intent remains underspecified, externalizing intent, interfaces, constraints,
and acceptance criteria may reduce ambiguity before it spreads across generated
code and parallel contributors~\cite{github2025spec,thoughtworks202Xspec}.

This interpretation remains compatible with Agile values. It argues for a
clearer object of iteration, not for a frozen plan: specifications can change,
but revisions should be explicit, reviewable, and connected to tasks and
implementation. For educators, the implication is to assess shared
architectural intent, API contracts, cross-component dependencies, design
revision rationales, and verifiable acceptance criteria before and during
GenAI-assisted implementation. Future classroom studies should pair process
metrics with concept assessments and structured reflections to distinguish
productive integration struggle from delayed recognition of foundational gaps.

\subsection{Scope of the interpretation}

The discussion is bounded by the evidence design: RQ1 uses repeated cohort
summaries, RQ2 combines repository and coverage-limited transcript evidence,
and RQ3 uses a 14 team-semester panel. These limits support
mechanism-oriented hypothesis generation, not causal estimates of GenAI,
planning, or SDD effects.
